% Capitulo: Por que o ágil de prateleira não encaixa em dados
\chapter{Por que o ágil de prateleira não encaixa em dados}
\label{cap:agil-prateleira}

Muitos times de dados já tentaram adotar o Scrum exatamente como descrito no guia: iterações
de duração fixa, estimativas em pontos, compromisso com um conjunto de histórias no começo de
cada ciclo. E muitos descobriram que algo não encaixa. As histórias mudam de tamanho no meio do
caminho, metade do time é interrompida por demandas urgentes, e a revisão no fim do ciclo
mostra mais trabalho inacabado do que entregue. Este capítulo explica por que isso acontece e
o que vale aproveitar de cada método. Não é uma crítica ao ágil, e sim ao ágil aplicado sem
adaptação.

\section{Descoberta e incerteza}
\label{sec:descoberta}

O ágil nasceu no desenvolvimento de software, e o manifesto que lhe deu nome valoriza as
pessoas, o software funcionando, a colaboração com o cliente e a resposta à mudança
\parencite{agilemanifesto2001}. Nada disso é estranho ao trabalho com dados. O que é diferente
é a natureza da incerteza.

Num projeto de software típico, a maior parte da incerteza está em como construir. Num projeto
de dados, uma parte grande está nos próprios dados, e só aparece quando alguém olha para eles.
Uma história como \enquote{criar o indicador de clientes ativos por região} parece pequena até
que se descobre que a região não está preenchida para um terço dos clientes, que há dois
cadastros de região que não batem e que a área comercial e a financeira definem
\enquote{ativo} de formas diferentes. A tarefa não ficou maior por erro de estimativa; ela
revelou o seu tamanho real quando os dados foram abertos.

A literatura sobre estimativa de software descreve essa incerteza como um cone: no começo de
um projeto, as estimativas podem errar por um fator grande para mais ou para menos, e a margem
só se estreita à medida que o trabalho avança e as incógnitas são resolvidas
\parencite{mcconnell2006}. Em dados, o cone é mais largo e se estreita mais tarde, porque
parte das incógnitas só se resolve olhando os dados. Um método que exige compromisso com o
escopo no começo de cada ciclo, com base em estimativas feitas antes de olhar os dados, está
pedindo ao time o que ele não tem como dar.

A consequência prática é separar descoberta de entrega. Antes de comprometer uma história de
construção, uma investigação curta, com prazo fixo, olha os dados e responde às perguntas que
definem o tamanho do trabalho. Só então a história de construção é refinada e estimada. O
capítulo~\ref{cap:scrumban} volta a essa separação.

\section{Dependências e sustentação}
\label{sec:dependencias}

Duas outras características do trabalho com dados atrapalham o ágil de prateleira. A primeira
são as dependências externas. Um time de dados depende das áreas donas dos sistemas de origem
para liberar acessos, explicar campos, corrigir cadastros e avisar mudanças, e depende das
áreas de negócio para definir regras e validar resultados. Nenhuma dessas pessoas está no time,
e nenhuma segue o ritmo das iterações dele. Uma história pode ficar parada dias esperando uma
resposta, sem que o time possa fazer nada.

A segunda é a sustentação. Um time de dados mantém em produção tudo o que já entregou:
pipelines, modelos, painéis. Quando uma carga falha ou um número sai estranho, a correção não
pode esperar o fim da iteração, porque alguém está decidindo com aquele número agora. Esse
trabalho não planejado chega todos os dias, em quantidade variável, e consome uma parte da
capacidade do time que nenhum planejamento consegue prever com precisão.

Num método que mede o sucesso pela entrega do que foi comprometido no começo do ciclo, as
dependências e a sustentação fazem o time parecer pior do que é. O planejamento fica sempre
otimista, a revisão mostra sempre trabalho inacabado, e o time passa a desconfiar do próprio
processo.

\rascunho{
  Relato: uma experiência com Scrum aplicado ao pé da letra num time de dados, e o que deu
  errado (histórias que mudavam de tamanho, interrupções de sustentação, revisões frustrantes).
}

\section{O que aproveitar de cada método}
\label{sec:aproveitar}

A saída não é abandonar o ágil, e sim combinar o que cada método tem de melhor para o trabalho
com dados. O Scrum e o Kanban, os dois mais usados, têm forças complementares.

Do Scrum, vale aproveitar as cadências \parencite{scrumguide2020}. Um planejamento periódico
obriga o time e o negócio a discutir prioridades com regularidade. A revisão, em que o
trabalho pronto é mostrado a quem pediu, mantém o negócio perto e permite corrigir o rumo cedo.
A retrospectiva, em que o time discute o próprio processo, é a forma de melhorar
continuamente. E a definição de pronto, um critério explícito do que precisa estar feito para
um item ser considerado entregue, evita a entrega que parece terminada e não está.

Do Kanban, vale aproveitar a gestão do fluxo \parencite{anderson2010}. Visualizar todo o
trabalho num quadro, inclusive a sustentação e as dependências, mostra onde as coisas estão
paradas. Limitar o trabalho em andamento força o time a terminar antes de começar, o que reduz
o tempo de cada item no sistema. Tornar as regras explícitas (o que entra em cada coluna,
quando um item está pronto para a próxima etapa) reduz discussões e mal-entendidos. E medir o
fluxo, com indicadores como o tempo que cada item leva do começo ao fim, dá uma base objetiva
para previsões, em vez de estimativas feitas no escuro.

A combinação dos dois, às vezes chamada de Scrumban, foi proposta originalmente como uma forma
de times acostumados ao Scrum evoluírem em direção a um sistema de fluxo \parencite{ladas2009}.
Para dados, ela resolve boa parte dos problemas descritos neste capítulo: as cadências mantêm o
negócio por perto, o fluxo acomoda a descoberta e a sustentação, e os limites de trabalho em
andamento protegem o time de começar tudo e terminar nada. O próximo capítulo descreve como
esse modelo pode ser montado na prática.

\begin{leituras}
  \leitura{anderson2010}{A referência do método Kanban para trabalho intelectual, com as
    práticas de visualizar, limitar o trabalho em andamento e gerir o fluxo.}
  \leitura{scrumguide2020}{O guia oficial do Scrum, curto e direto. Vale reler para separar
    o que o Scrum exige do que foi acrescentado pela prática.}
  \leitura{mcconnell2006}{Para entender por que estimativas no começo de um projeto erram
    tanto, e como lidar com a incerteza em vez de ignorá-la.}
\end{leituras}
